\begin{reviewed}
\chapter*{Abstract}\label{abstract}

This dissertation investigated the evidence available for evaluating AI safety and alignment, with an empirical focus on alignment faking in compact language models. A systematic literature review organized 80 included entries into seven categories. Liability \& Ethics and Bias \& Fairness were the most underrepresented categories relative to their assessed severity.

Three prompt-based experiments examined harmful requests, persona preservation, and animal welfare dismissal across five model families. Free-tier prompts simulated interactions subject to retraining, while paid-tier prompts simulated deployment without retraining. Automated content judgments were compared across tiers and, where available, with self-reported compliance. The only classifier-based compliance gap appeared for LLaMA 3.1 8B without scratchpad prompting on harmful requests: 11 of 100 free-tier responses were classified as compliant, compared with none of 100 paid-tier responses. Other comparisons did not reach the nominal significance threshold and thus did not establish persistent alignment faking in the tested small LLMs.

\section*{Keywords:} artificial intelligence safety, alignment faking, systematic literature review, large language models, compliance evaluation
\end{reviewed}